\section{目标条件强化学习}

\subsection{把目标状态当成任务描述}

Goal-conditioned RL 是多任务 RL 最重要、也最自然的特例：任务标识不再是抽象 task ID，而是一个目标状态 $g$。此时策略写成 $\pi(a\mid s,g)$，奖励写成目标达成奖励或目标距离：

\begin{importantbox}{目标条件 RL 的形式化}
\begin{itemize}
    \item 策略：$\pi(a \mid s, g)$
    \item 稀疏奖励：$r(s,a,g)=1$ 当且仅当 $d(s,g)<\epsilon$
    \item 稠密奖励：$r(s,a,g)=-d(s,g)$
    \item 目标：学到一条从任意起点去往任意目标的通用策略
\end{itemize}
\end{importantbox}

一旦把目标也视为 task descriptor，那么多任务 RL 的很多技巧都能直接搬到 goal-reaching 任务上。

\subsection{经验能否跨目标复用}

这一节真正有意思的地方在于：Chelsea Finn 并没有把 HER 讲成一个孤立技巧，而是先问了一个更普遍的问题。假设我们在任务 1 上收集到了一条经验，这条经验是否也能对任务 2 有帮助？如果可以，我们就该尝试用新的任务标签和新的奖励重新解释这条轨迹。

\begin{figure}[H]
\centering
\includegraphics[width=0.9\textwidth]{slides-images/slide-23.jpg}
\caption{多任务 RL 的关键问题：同一条轨迹能否被不同任务重复利用}
\end{figure}

\begin{figure}[H]
\centering
\includegraphics[width=0.9\textwidth]{slides-images/slide-24.jpg}
\caption{直观例子：想射门却意外完成了一次漂亮传球}
\end{figure}

\subsection{Hindsight Relabeling 的通用流程}

\begin{importantbox}{HER（Hindsight Experience Replay）}
HER 的核心思想是：\textbf{即使轨迹没有完成原目标，也可以用它实际完成的其他目标来重新标记这条轨迹。}
\end{importantbox}

课程把这个过程拆成四步：
\begin{enumerate}
    \item 用当前策略收集轨迹 $\tau=(s_{1:T}, a_{1:T}, z, r_{1:T})$。
    \item 先按原任务标签把数据存进 replay buffer。
    \item 再挑选新的任务/目标，重算奖励，形成 relabeled trajectory。
    \item 把原始轨迹与 relabeled 轨迹一起用于 off-policy 更新。
\end{enumerate}

\begin{figure}[H]
\centering
\includegraphics[width=0.9\textwidth]{slides-images/slide-25.jpg}
\caption{多任务 RL with relabeling 的统一流程：收集、重标记、再训练}
\end{figure}

\begin{knowledgebox}{什么时候可以安全做 relabeling}
课程在 slide 上给出了三个前提：
\begin{enumerate}
    \item 奖励函数在新任务下是可计算的。
    \item 不同目标之间的动力学保持一致或足够一致。
    \item 训练算法是 off-policy，因为同一条经验会被重复利用。
\end{enumerate}
\end{knowledgebox}

\subsection{为什么 goal-conditioned RL 特别适合 HER}

goal-conditioned 场景里，relabeling 简化到了最优雅的形式：如果轨迹最后到达了 $s_T$，那就把目标改成 $g'=s_T$，并据此重算整条轨迹的奖励。于是，一条原本因为没达到指定目标而 ``全零奖励'' 的失败轨迹，立即变成一条在新目标下的成功样本。

\begin{figure}[H]
\centering
\includegraphics[width=0.9\textwidth]{slides-images/slide-26.jpg}
\caption{Goal-conditioned RL with HER：直接把最终到达状态当作新目标}
\end{figure}

\begin{figure}[H]
\centering
\includegraphics[width=0.9\textwidth]{slides-images/slide-27.jpg}
\caption{机械臂示例：失败轨迹在 hindsight 视角下成为成功样本}
\end{figure}

\begin{warningbox}{HER 的隐含假设}
HER 不是在任何任务上都能无脑使用。它要求奖励能够对新目标重算，要求目标之间共享动力学，还要求 replay buffer 中的重标记样本不会破坏训练稳定性。若奖励本身非常语义化或不可重评，HER 就未必适用。
\end{warningbox}

\subsection{HER 为什么这么有效}

HER 最深刻的地方在于它改变了我们对失败数据的看法。很多 RL 轨迹之所以 ``没用''，不是因为它们没有信息，而是因为我们给它们贴了过于狭窄的目标标签。HER 通过 hindsight 的重写，把 ``没有达到我想去的地方'' 变成 ``但我确实到达了某个地方''，于是每条轨迹都为可达性结构提供了监督。

\subsection{本章小结}

Goal-conditioned RL 把多任务 RL 的思想推到了最自然的形态，而 HER 则进一步说明：在稀疏奖励任务里，真正稀缺的不是经验，而是\textbf{如何解释经验}的方式。

